
\section{Method} \label{section:method}

In this section, we present Decision Transformer, which models trajectories autoregressively with minimal modification to the transformer architecture, as summarized in Figure \ref{fig:main_fig} and Algorithm \ref{alg:decisiontransformer}.

\textbf{Trajectory representation.}
The key desiderata in our choice of trajectory representation are
that it should enable transformers to learn meaningful patterns and we should be able to conditionally generate actions at test time.
It is nontrivial to model rewards since we would like the model to generate actions based on \emph{future} desired returns, rather than past rewards.
As a result, instead of feeding the rewards directly, we feed the model with the returns-to-go $\widehat{R}_t = \sum_{t'=t}^T r_{t'}$.
This leads to the following trajectory representation which is amenable to autoregressive training and generation:
\begin{equation}
    \tau = \left(\widehat{R}_1, s_1, a_1, \widehat{R}_2, s_2, a_2, \dots, \widehat{R}_T, s_T, a_T\right).
\end{equation}

At test time, we can specify the desired performance (e.g. 1 for success or 0 for failure), as well as the environment starting state, as the conditioning information to initiate generation.
After executing the generated action for the current state, we decrement the target return by the achieved reward and repeat until episode termination.

\textbf{Architecture.}
We feed the last $K$ timesteps into Decision Transformer, for a total of $3K$ tokens (one for each modality: return-to-go, state, or action).
To obtain token embeddings, we learn a linear layer for each modality, which projects raw inputs to the embedding dimension, followed by layer normalization~\citep{ba2016layernorm}. 
For environments with visual inputs, the state is fed into a convolutional encoder instead of a linear layer.
Additionally, an embedding for each timestep is learned and added to each token -- note this is different than the standard positional embedding used by transformers, as one timestep corresponds to three tokens.
The tokens are then processed by a GPT \citep{radford2018gpt} model, which predicts future action tokens via autoregressive modeling.

\textbf{Training.}
We are given a dataset of offline trajectories.
We sample minibatches of sequence length $K$ from the dataset.
The prediction head corresponding to the input token $s_t$ is trained to predict $a_t$ -- either with cross-entropy loss for discrete actions or mean-squared error for continuous actions -- and the losses for each timestep are averaged.
We did not find predicting the states or returns-to-go to improve performance, although
it is easily permissible within our framework (as shown in Section \ref{sec:key_to_door}) and would be an interesting study for future work.

\begin{algorithm}[hb]
\caption{Decision Transformer Pseudocode (for continuous actions)}
\label{alg:decisiontransformer}
\definecolor{codeblue}{rgb}{0.28125,0.46875,0.8125}
\lstset{
    basicstyle=\fontsize{9pt}{9pt}\ttfamily\bfseries,
    commentstyle=\fontsize{9pt}{9pt}\color{codeblue},
    keywordstyle=
}
\begin{lstlisting}[language=python]
# R, s, a, t: returns-to-go, states, actions, or timesteps
# transformer: transformer with causal masking (GPT)
# embed_s, embed_a, embed_R: linear embedding layers
# embed_t: learned episode positional embedding
# pred_a: linear action prediction layer

# main model
def DecisionTransformer(R, s, a, t):
    # compute embeddings for tokens
    pos_embedding = embed_t(t)  # per-timestep (note: not per-token)
    s_embedding = embed_s(s) + pos_embedding
    a_embedding = embed_a(a) + pos_embedding
    R_embedding = embed_R(R) + pos_embedding
    
    # interleave tokens as (R_1, s_1, a_1, ..., R_K, s_K)
    input_embeds = stack(R_embedding, s_embedding, a_embedding)
    
    # use transformer to get hidden states
    hidden_states = transformer(input_embeds=input_embeds)
    
    # select hidden states for action prediction tokens
    a_hidden = unstack(hidden_states).actions
    
    # predict action
    return pred_a(a_hidden)

# training loop
for (R, s, a, t) in dataloader:  # dims: (batch_size, K, dim)
    a_preds = DecisionTransformer(R, s, a, t)
    loss = mean((a_preds - a)**2)  # L2 loss for continuous actions
    optimizer.zero_grad(); loss.backward(); optimizer.step()

# evaluation loop
target_return = 1  # for instance, expert-level return
R, s, a, t, done = [target_return], [env.reset()], [], [1], False
while not done:  # autoregressive generation/sampling
    # sample next action
    action = DecisionTransformer(R, s, a, t)[-1]  # for cts actions
    new_s, r, done, _ = env.step(action)
    
    # append new tokens to sequence
    R = R + [R[-1] - r]  # decrement returns-to-go with reward
    s, a, t = s + [new_s], a + [action], t + [len(R)]
    R, s, a, t = R[-K:], ...  # only keep context length of K
\end{lstlisting}
\end{algorithm}
